\section{From activity packing to a power saving}
\label{sec:recurrence}

We now turn the uniform activity-packing estimate into a fixed power
improvement. The scale is chosen once and then used at every size in
an induction. Throughout this section, a sweep is the ordering by
increasing values of $y-sx$ as $s$ runs from $-\infty$ to $+\infty$.
Its rank-$k$ marked set consists of the $k$ smallest values. The
\emph{activity} of an interval of sweep steps is the number of distinct
points whose membership in that set changes during the interval.

\begin{proposition}\label{prop:packing-to-power}
Suppose that, for every $\eta>0$, there is $R_\eta<\infty$ such that
for every generic planar configuration of $N$ points, every rank, and
every integer $1\le b\le N$ with $N/b\ge R_\eta$, the number of
step-disjoint intervals of activity at least $b$ is at most
$\eta(N/b)^{4/3}$. Then there are absolute constants $\varepsilon>0$
and $C$ such that every such rank sweep has at most
$CN^{4/3-\varepsilon}$ switches.
\end{proposition}

\begin{proof}
Write $a=4/3$, set $\eta=(4\cdot2^a)^{-1}$, and choose an integer
$q\ge\max\{4,R_\eta\}$. This integer remains fixed. For $N\ge4q$,
put
\[
 d=\lfloor N/q\rfloor,\qquad b=d-1.
\]
Then $N/(2q)\le b<N/q$, so the packing hypothesis applies and
$N/b\le2q$.

Partition the sweep greedily, ending a block when its activity first
reaches $b$, and include any final remainder. Immediately before the
last step of a completed block, the activity is at most $b-1$.
A step changes at most two memberships, so the block has at most
$b+1=d$ active points. The remainder has fewer than $b$ active points.
The completed blocks are step-disjoint, and hence their number is at
most
\[
 \eta(N/b)^a\le\eta(2q)^a=q^a/4.
\]
Including the one possible remainder gives at most $q^a/2$ blocks,
since $q^a\ge4$. Every block uses at most $N/q$ active points.

Let $A$ be the active points of one block and let $B$ be the points
outside $A$ that are marked at its initial state. Membership of every
point outside $A$ is constant during the block. Therefore restriction
to $A$ is a sweep at the fixed rank $r=k-|B|$: at every state, the
marked points of $A$ are exactly its $r$ lowest points. This identity
also proves that the switches in the block are exactly the switches
of the restricted rank-$r$ sweep over that slope interval. Their
number is bounded by the complete-sweep complexity on $|A|$ points.

Let $F(n)$ be the largest switch count among all generic configurations
of exactly $n$ points and all ranks, with $F(0)=F(1)=0$. It is finite,
since $F(n)\le\binom n2$. Set
\[
 \varepsilon=\frac{\log2}{2\log q},\qquad
 \beta=\frac43-\varepsilon>0.
\]
Choose $C$ so that $F(n)\le Cn^\beta$ for $1\le n<4q$. Suppose this
bound holds at every size smaller than $N\ge4q$. If $t_j\le N/q$
are the active-subset sizes of the blocks, the restriction argument
and the induction hypothesis give
\[
 F(N)\le\sum_j F(t_j)
 \le C\sum_j t_j^\beta
 \le C\frac{q^a}{2}\left(\frac Nq\right)^\beta
 =\frac{C}{\sqrt2}N^\beta
 \le CN^\beta.
\]
Here every nonempty subset has size less than $N$, while empty subsets
contribute zero. Thus induction proves the assertion for all sizes
and ranks. The selection of $q$ from $R_\eta$ is nonquantitative;
this argument does not provide a numerical value of $\varepsilon$.
\end{proof}

Proposition~\ref{prop:packing-to-power} and Theorem~\ref{thm:activity-packing}
prove Theorem~\ref{thm:all-ranks}. It remains to transfer its middle-rank
case to the configurations of Theorem~\ref{thm:halving-power}.

For a generic configuration of $N=2m$ points, a pair is halving
precisely when its critical slope ties positions $m$ and $m+1$:
there are then $m-1$ points on each side of its line. The two tied
points exchange order, so halving pairs are exactly the middle-rank
switches. Each unordered pair has one critical slope in this sweep.
The perturbation argument in Section~\ref{sec:finite} preserves all
halving pairs of an arbitrary configuration with no three collinear. Consequently
Proposition~\ref{prop:packing-to-power}, applied to Theorem~\ref{thm:activity-packing}, gives the asserted bound for every even
size and every such configuration.

\subsection{Planar \texorpdfstring{$k$}{k}-sets and straight-line levels}
\label{sec:k-sensitive-proof}

We finish by deriving the rank-sensitive consequence stated in the
introduction. The local ranks created by a shallow cutting can vary from
cell to cell, which is why the uniform sweep theorem was needed.

\begin{proof}[Proof of Corollary~\ref{cor:k-sensitive}]
Take $\varepsilon$ from Theorem~\ref{thm:all-ranks} and set
$\varepsilon_0=\min\{\varepsilon,1/6\}$ and
$\beta=4/3-\varepsilon_0>1$. The theorem remains valid with exponent
$\beta$ in place of $4/3-\varepsilon$.

Dualize a point $(a,b)$ to the line $y=b-as$. For a generic dual point
set $P$, let $V_j$ be the arrangement vertices with exactly $j$ lines
strictly below them. Then $|V_j|=S_{j+1}(P)$: at such a vertex the lines
in positions $j+1$ and $j+2$ exchange order. The vertices of
$\Lambda_j$ are $V_{j-1}\cup V_j$, with out-of-range sets empty
\cite[Section~1, p.~2]{AgarwalAronovChanSharir1998}. Thus
Theorem~\ref{thm:all-ranks} bounds every level of $M$ dual lines by
$O(M^\beta)$ vertices, uniformly in its rank; its number of edge pieces
is one more than its number of vertices. For any line arrangement in
the corollary, the dual point set has no three collinear and distinct
first coordinates. The perturbation in Section~\ref{sec:finite} can
also preserve the order of those coordinates, so its preservation of
all triple-orientation signs preserves every dual vertex's number of
lines below it. The uniform $O(M^\beta)$ level bound therefore holds
for every such general-position line arrangement.

Apply Matou\v{s}ek's shallow-cutting lemma in the form stated by
Agarwal, Aronov, Chan, and Sharir
\cite[Lemma~3.2 and the following reduction, pp.~10--11]{AgarwalAronovChanSharir1998}.
In dimension two, take $r=n/(k+1)\ge1$ and
$q=kr/n+1<2$. The cutting has $O(rq)=O(n/(k+1))$ cells, each crossed
by at most $n/r=k+1$ lines. Within the interior of a cell, the
noncrossing lines give a fixed rank offset, so the part of the global
level has one fixed local rank in the arrangement of crossing lines.
This local rank need not
equal $k$, which is why the all-ranks bound is used. Boundary pieces
cost $O(k+1)$ per cell: each crossing line meets the simplex boundary
at most twice. An input line not crossing the interior can contribute
only along a side or at a vertex. General position allows at most one
such line per side and at most two through each vertex; the already
counted intersections split their contributions into $O(k+1)$ pieces
in total. Since $\beta>1$, these terms are absorbed by the local
$O((k+1)^\beta)$ bound. The cited shallow-cutting reduction now gives
\[
 \operatorname{comp}\bigl(\Lambda_k(\mathcal L)\bigr)
 =O\bigl(r(n/r)^\beta\bigr)
 =O\bigl(n(k+1)^{1/3-\varepsilon_0}\bigr).
\]

For a generic point set $Q$, every $k$-set consists of the $k$ smallest
values of some linear functional. Normal directions with positive
second coordinate give the bottom-$k$ sweep, while those with negative
second coordinate give complements of the bottom-$(n-k)$ sweep.
The horizontal directions are included by limits, since first
coordinates are distinct. Hence the number of $k$-sets is at most
$S_k(Q)+S_{n-k}(Q)+2$. In the dual arrangement,
$S_k(Q)=|V_{k-1}|$; vertical reflection sends the vertices counted by
$S_{n-k}(Q)=|V_{n-k-1}|$ to $V_{k-1}$ of the reflected arrangement.
Each of these two vertex sets
lies on a $k$-th level, so the level bound just proved applies twice.
Finally, for any fixed point set, choose one strict separating line for
each of its finitely many $k$-sets. All these strict separations survive
a sufficiently small generic perturbation. Its original number of
$k$-sets is therefore no larger than the generic count.
\end{proof}
